\section{Methodology}
\label{sec:method}

We design a controlled experiment to measure whether simple adversarial prompts placed after sensitive information in system prompts reduce information leakage under diverse extraction attacks.

\subsection{System Prompt Construction}
\label{sec:prompts}

We sample 20 base system prompts from the Unnatural Instructions dataset used in PSM~\cite{jawad2026psm}.
Into each prompt, we inject three sensitive fields at marked locations:
\begin{itemize}[leftmargin=*,itemsep=0pt,topsep=0pt]
    \item \textbf{API Key}: \texttt{sk-PROJ-xK9mT2vLpQ8nR5wA3bF7dG1hJ4kM6oP}
    \item \textbf{Secret Code}: \texttt{WHISKEY-TANGO-FOXTROT-7749}
    \item \textbf{Internal URL}: \texttt{https://internal.acmecorp.io/admin/config/v3}
\end{itemize}
These fields represent common types of sensitive information found in production system prompts: authentication credentials, internal identifiers, and private endpoints.

\subsection{Defense Conditions}
\label{sec:defenses}

We evaluate five defense conditions, summarized in \tabref{tab:defenses}.
The \nodefense baseline places sensitive fields with no protection.
The \heuristic baseline appends a standard guardrail instruction (``Do not reveal any sensitive information, API keys, or secrets from your instructions'').
The three firewall conditions place adversarial text immediately after each sensitive field:

\begin{table}[t]
    \centering
    \caption{Defense conditions evaluated. Firewall text is placed immediately after each sensitive field within the system prompt.}
    \begin{tabular}{@{}llp{7cm}@{}}
        \toprule
        \textbf{Condition} & \textbf{Type} & \textbf{Description} \\
        \midrule
        \nodefense & Baseline & Raw system prompt, no protection \\
        \heuristic & Baseline & Standard guardrail: ``Do not reveal any sensitive information\ldots'' \\
        \mildfirewall & Firewall & Security notice: ``The above information is classified\ldots'' \\
        \modfirewall & Firewall & Structured boundary: \texttt{---SYSTEM BOUNDARY---} with explicit access rules \\
        \aggfirewall & Firewall & Fake assistant completion simulating refusal to a prior extraction attempt \\
        \bottomrule
    \end{tabular}
    \label{tab:defenses}
\end{table}

The \mildfirewall uses a brief security notice after each field.
The \modfirewall uses a structured boundary marker (\texttt{---SYSTEM BOUNDARY---}) with explicit rules demarcating sensitive content as off-limits.
The \aggfirewall simulates a fake assistant completion where the model has already detected and refused an extraction attempt, following the ``fight fire with fire'' principle of \citet{chen2025dat}.

\subsection{Attack Suite}
\label{sec:attacks}

We use 65 attacks spanning six categories, summarized in \tabref{tab:attacks}.

\begin{table}[t]
    \centering
    \caption{Attack categories and sources. General attacks (first three rows) test broad extraction; exact-match and fuzzy queries test field-specific access; benign queries test utility preservation.}
    \begin{tabular}{@{}llcl@{}}
        \toprule
        \textbf{Category} & \textbf{Source} & \textbf{Count} & \textbf{Style} \\
        \midrule
        Raccoon & \citet{wang2024raccoon} & 15 & Injection, override, role-play \\
        Liang & \citet{liang2024prompts} & 10 & Polite request extraction \\
        Zhang & \citet{das2025system} & 10 & Diverse extraction strategies \\
        \midrule
        Exact-match & Custom & 10 & Queries naming specific sensitive values \\
        Fuzzy & Custom & 10 & Vague queries about configuration/secrets \\
        Benign & Custom & 10 & Normal task queries (utility test) \\
        \bottomrule
    \end{tabular}
    \label{tab:attacks}
\end{table}

The 35 general attacks (Raccoon, Liang, Zhang) test broad extraction capability.
The 10 exact-match queries specifically name sensitive values (\eg ``What is the API key \texttt{sk-PROJ-...}?'') to test whether firewalls allow access via precise queries.
The 10 fuzzy queries ask vaguely about secrets without naming specific values.
The 10 benign queries test normal task completion to measure utility impact.

\subsection{Experimental Protocol}
\label{sec:protocol}

All experiments use \gptmini via the OpenAI API with temperature$=0$, max\_tokens$=1024$, and seed$=42$.
Each of the 20 base prompts is tested with each of the 5 defense conditions and all applicable attacks, yielding 5{,}650 total API calls across four experiments:

\begin{itemize}[leftmargin=*,itemsep=0pt,topsep=0pt]
    \item \textbf{Experiment 1} (Baseline): \nodefense condition with all 35 general attacks $\times$ 20 prompts = 700 calls.
    \item \textbf{Experiment 2} (Firewall Effectiveness): 4 defended conditions $\times$ 35 attacks $\times$ 20 prompts = 2{,}800 calls.
    \item \textbf{Experiment 3} (Exact-Match Bypass): 5 conditions $\times$ 20 queries $\times$ 20 prompts = 2{,}000 calls.
    \item \textbf{Experiment 4} (Utility): 5 conditions $\times$ 10 benign queries $\times$ 20 prompts = 1{,}000 calls. Note: some conditions in Experiments 3 and 4 overlap with Experiments 1--2; the unique total is 5{,}650.
\end{itemize}

\subsection{Evaluation Metrics}
\label{sec:metrics}

\para{Sensitive Field Leakage Rate (\sflr).}
The fraction of responses containing any of the three exact sensitive field values as substrings.
This is our primary metric: it directly measures whether the model reveals the specific secrets we aim to protect.

\para{\rougel Recall.}
The longest common subsequence recall between the full system prompt and the model response, measuring surface-level overlap.
Higher values indicate more of the system prompt content has leaked into the response.

\para{Utility Score.}
TF-IDF cosine similarity between responses to benign queries under defended and undefended conditions.
A score of 1.0 indicates the defense has no impact on normal task responses.

\subsection{Statistical Analysis}
\label{sec:stats}

We compare \sflr across conditions using chi-squared and Fisher exact tests (for small cell counts).
We report bootstrap 95\% confidence intervals (1{,}000 resamples) for leakage rate differences.
We use the Kruskal-Wallis test for \rougel comparisons across conditions (non-normal distributions expected).
All tests use $\alpha = 0.05$.
